\chapter{Zadatak 4 --- što je točno zadano}
\label{ch:z4-zadatak}

\section{Izvorni tekst, doslovno}
\label{sec:z4-izvorni}

\begin{zadatak}{Zadatak 4 (Robusna stabilizacija)}
\skenzad[0.96]{z4-original}
\end{zadatak}

\subsection{Prijepis}

\begin{zadatak}{Zadatak 4 --- prijepis}
\textbf{Zadatak 4. (Robusna stabilizacija)} Promotrimo sljedeći vremenski
promjenjiv sustav
\begin{equation}
\dot{x}(t) = A(t)x(t) + Bu(t) \tag{3}
\end{equation}
gdje je $x(t)$ varijabla stanja, a $u(t)$ je ulazni (upravljački) signal. Matrica
$A(t)$ je vremenski promjenjiva ($t$ označava vrijeme) i to prema sljedećem
pravilu:
\begin{equation}
A(t) = \begin{cases}
A_{\mathrm{parno}} & \text{kad je } \lfloor t \rfloor \text{ paran broj} \\
A_{\mathrm{neparno}} & \text{kad je } \lfloor t \rfloor \text{ neparan broj}
\end{cases} \tag{4}
\end{equation}
gdje $\lfloor t \rfloor$ označava najveći cijeli broj koji je manji ili jednak
$t$. Imamo:
\begin{equation}
A_{\mathrm{parno}} = \begin{bmatrix} -1 & -\tfrac{1}{6}\pi \\ \tfrac{3}{2}\pi & -1 \end{bmatrix},
\qquad
A_{\mathrm{neparno}} = \begin{bmatrix} -1 & -\tfrac{3}{2}\pi \\ \tfrac{1}{6}\pi & -1 \end{bmatrix}. \tag{5}
\end{equation}
Ovdje se u biti radi o hibridnom sustavu sa ,,skokovitom`` promjenom dinamike
(\emph{switching}) u diskretnim vremenskim trenutcima. Između svaka dva trenutka
u kojima se odvija ova skokovita promjena, sustav je linearan vremenski
invarijantan. Kažemo da ovaj sustav ima dva dinamička moda.

\begin{enumerate}[label=\alph*),leftmargin=*]
  \item Jesu li matrice $A_{\mathrm{parno}}$ i $A_{\mathrm{neparno}}$ stabilne
        matrice (Hurwitzove matrice)? Odnosno, jesu li oba dinamička moda
        stabilna? \emph{(Napomena: ovo možete provjeriti izračunom svojstvenih
        vrijednosti ovih matrica.)}
  \item Simulirajte vremenski promjenjiv sustav (3) s početnim uvjetom
        $x(0) = \begin{bsmallmatrix} 1 \\ 0 \end{bsmallmatrix}$, prikažite $x(t)$
        (graf kako se varijable stanja mijenjaju s vremenom) i prikažite
        trajektoriju sustava u prostoru stanja. Je li ovaj sustav stabilan?
  \item U sustavu (3), neka je matrica
        $B = \begin{bsmallmatrix} 1 \\ 0 \end{bsmallmatrix}$. Dizajnirajte
        statički regulator $u(t) = Kx(t)$ takav da je zatvoreni krug
        \begin{equation}
          \dot{x}(t) = \bigl(A(t) + BK\bigr)x(t) \tag{6}
        \end{equation}
        eksponencijalno stabilan, gdje je $A(t)$ zadana s (4). Simulirajte ovaj
        vremenski promjenjiv zatvorni sustav s početnim uvjetom
        $x(0) = \begin{bsmallmatrix} 1 \\ 0 \end{bsmallmatrix}$, prikažite $x(t)$
        i prikažite trajektoriju sustava u prostoru stanja.
\end{enumerate}
\end{zadatak}

\section{Što u zadatku \emph{ne} piše}
\label{sec:z4-nije-receno}

\begin{longtable}{@{}p{0.36\linewidth} p{0.60\linewidth}@{}}
\toprule
\textbf{Što nije rečeno} & \textbf{Kako je protumačeno i zašto} \\
\midrule
\endfirsthead
\toprule
\textbf{Što nije rečeno} & \textbf{Kako je protumačeno i zašto} \\
\midrule
\endhead
Koliki je $B$ u podzadacima a) i b). & Zadaje se tek u c). U a) i b) je
$u \equiv 0$, pa $B$ ni ne treba --- promatra se autonomni sustav
$\dot{x} = A(t)x$. \\[6pt]
\textbf{Kojom metodom} projektirati $K$. & Ne kaže se. Tražena je samo
eksponencijalna stabilnost. Odabrana je zajednička kvadratna Ljapunovljeva
funkcija i sinteza preko LMI-ja, po Predavanju 7 --- jedina metoda koja
jamči stabilnost za \emph{svaki} zakon prespajanja, ne samo za zadani. \\[6pt]
Koliko dugo simulirati. & Nije zadano. Uzeto je $t \in [0, 20]$ s: dovoljno da
se u otvorenom krugu vidi rast, a u zatvorenom da norma padne ispod
$10^{-11}$. \\[6pt]
Kojim rješavačem. & Nije zadano. Simulacija po segmentima s \texttt{ode45}
(unutar segmenta sustav je LTI), LMI s YALMIP-om i SDPT3. \\[6pt]
Kako riješiti stroge nejednakosti numerički. & Rješavači rade s nestrogim
nejednakostima. Uvedene su normalizacije $\mathcal{Q} \psd I$ i
$M_i \nsd -I$; opravdanje je homogenost sustava
(\S\ref{sec:z4-normalizacija}). \\[6pt]
Koji $K$ odabrati među beskonačno mnogo dopustivih. & Sustav LMI-ja je problem
\emph{dopustivosti}, ne optimizacije --- rješenja ima beskonačno mnogo. U
seminaru se bira ono s najmanjom $\norm{Y}_2$. \\[6pt]
Treba li dokazati nestabilnost u b). & Traži se samo ,,je li ovaj sustav
stabilan``. Simulacija sugerira odgovor; u seminaru je dodan i dokaz preko
monodromijske matrice, jer simulacija sama po sebi nije dokaz. \\
\bottomrule
\end{longtable}

\begin{zamka}[Zamka: $\lfloor t \rfloor$ paran znači $t \in [0,1) \cup [2,3) \cup \dots$]
Prespajanje se događa na \emph{cjelobrojnim sekundama}, a mod vrijedi na
poluotvorenom intervalu. Za $t \in [0,1)$ je $\lfloor t \rfloor = 0$, paran, pa
vrijedi $A_{\mathrm{parno}}$; za $t \in [1,2)$ je $\lfloor t \rfloor = 1$, pa
vrijedi $A_{\mathrm{neparno}}$. Period prespajanja je dakle $2$ s, ne $1$ s --- to
je bitno za monodromijsku matricu u \S\ref{sec:z4-monodromija}.
\end{zamka}

\begin{zamka}[Zamka: ovo nije ,,nesigurnost``, nego \emph{prespajanje}]
Zadatak se zove ,,robusna stabilizacija``, ali $A(t)$ nije nepoznata konstanta
unutar nekog skupa --- ona se stvarno \emph{mijenja u vremenu}, po poznatom
zakonu. Razlika je ključna: kad bi $A$ bila nepoznata konstanta, dovoljno bi bilo
da svaka pojedina matrica bude Hurwitzova. Kod prespajanja to nije dovoljno, kao
što podzadatak b) i pokazuje.
\end{zamka}

\section{Karta poglavlja}
\label{sec:z4-karta}

\begin{longtable}{@{}p{0.07\linewidth} p{0.42\linewidth} p{0.45\linewidth}@{}}
\toprule
& \textbf{Podzadatak} & \textbf{Gdje je obrađen} \\
\midrule
\endhead
a) & Hurwitzovost modova & Poglavlje~\ref{ch:z4-hurwitz} \\
b) & Simulacija i (ne)stabilnost & Poglavlje~\ref{ch:z4-simulacija} \\
c) & Sinteza regulatora preko LMI-ja & Poglavlje~\ref{ch:z4-lmi} \\
--- & Provjera rješenja i simulacija zatvorenog kruga & Poglavlje~\ref{ch:z4-provjera} \\
--- & Pitanja, zamke i modelski odgovori & Poglavlje~\ref{ch:z4-obrana} \\
\bottomrule
\end{longtable}
